\section{Verification Approach}
\label{sec:approach}

\subsection{Methodology}

Our verification strategy combines three complementary techniques:

\begin{enumerate}[label=(\roman*)]
  \item \textbf{Mathematical proofs via Hoare logic.}
    Each function is verified using pre/postcondition reasoning with
    explicitly stated loop invariants. We prove initialization, maintenance,
    and termination for every loop. Termination arguments use integer-valued
    variant functions bounded below, establishing well-founded descent on $(\N, <)$.

  \item \textbf{Key lemmas connecting sorted arrays to the h-index.}
    We establish structural lemmas showing that for reverse-sorted arrays,
    the h-index coincides with the transition point where $a[h] \leq h$,
    reducing the general counting definition to a positional check.

  \item \textbf{Multiset reasoning for update correctness.}
    For the \texttt{update} function, we track the multiset of array values
    to establish that the update is equivalent to incrementing a citation count
    and re-sorting.
\end{enumerate}

\subsection{Notation and Conventions}

Throughout this document, we use the following conventions:

\begin{itemize}
  \item Arrays are 0-indexed: $a[0], a[1], \ldots, a[n-1]$.
  \item $a[i..j)$ denotes the subarray $a[i], a[i{+}1], \ldots, a[j{-}1]$.
  \item All variables are integers unless stated otherwise.
  \item $\lbag \cdot \rbag$ denotes multisets.
  \item We write $a \xrightarrow{\text{op}} a'$ to indicate that operation \texttt{op}
    transforms array~$a$ into array~$a'$.
  \item $|S|$ denotes the cardinality of set~$S$.
  \item Hoare triples $\{P\}\;\mathtt{S}\;\{Q\}$ assert partial correctness;
    total correctness (including termination) is denoted $[P]\;\mathtt{S}\;[Q]$.
\end{itemize}

\subsection{Challenge Functions}

The three C functions under verification are:

\begin{lstlisting}[style=cstyle, caption={Linear scan h-index computation}]
int compute(int a[], int n) {
    int h = 0;
    while (h < n && h < a[h]) h++;
    return h;
}
\end{lstlisting}

\begin{lstlisting}[style=cstyle, caption={Binary search h-index computation}]
int compute_opt(int a[], int n) {
    int lo = 0, hi = n;
    while (lo < hi) {
        int mid = lo + (hi - lo) / 2;
        if (a[mid] <= mid) hi = mid;
        else lo = mid + 1;
    }
    return lo;
}
\end{lstlisting}

\begin{lstlisting}[style=cstyle, caption={Incremental update with h-index maintenance}]
int update(int a[], int h, int i) {
    int x = a[i];
    int lo = 0, hi = i;
    while (lo < hi) {
        int mid = lo + (hi - lo) / 2;
        if (a[mid] == x) hi = mid;
        else lo = mid + 1;
    }
    a[lo]++;
    if (lo == h && a[lo] == h+1) return h+1;
    return h;
}
\end{lstlisting}
